\icode{N23L-MP-SI-L2-261004-002}
\nceeID[4]{2023辽宁}{6}
原子处于磁场中，某些能级会发生劈裂。某种原子能级劈裂前后的部分能级图如图所示，相应能级跃迁放出的光子分别设为\ding{172}\ding{173}\ding{174}\ding{175}。若用\ding{172}照射某金属表面时能发生光电效应，且逸出光电子的最大初动能为\(E_k\)，则\paren[A]

\begin{tikzpicture}[>=stealth,font=\scriptsize]
	\draw[dashed] (-2,0)--(5,0) (-2,-3)--(5,-3);
	\draw (-1.5,0)--node[above]{劈裂前}(.5,0) (2,0)--(4,0) (-1.5,-3)--(.5,-3) (2,-3)--(4,-3) (2,.2)--node[above]{劈裂后}(4,.2) (2,-.2)--(4,-.2);
	\draw[->] (-.5,0)--(-.5,-3)node[below]{\ding{172}};
	\draw[->] (3,0)--(3,-3)node[below]{\ding{174}};
	\draw[->] (2.3,-.2)--(2.3,-3)node[below]{\ding{173}};
	\draw[->] (3.7,.2)--(3.7,-3)node[below]{\ding{175}};
\end{tikzpicture}
\begin{choices}
	\item \ding{172}和\ding{174}的能量相等
	\item \ding{173}的频率大于\ding{175}的频率
	\item 用\ding{173}照射该金属一定能发生光电效应
	\item 用\ding{175}照射该金属逸出光电子的最大初动能小于\(E_k\)
\end{choices}
\begin{answer}
    A\qquad 提示：频率\(\nu_4 >\nu_1 =\nu_3 >\nu_2\)
\end{answer}